\input{../tex-inputs/latex-leseansicht-vorspann}

\section[Arthur und Louise Schnitzler an Gustav Schwarzkopf, 17. 4. 1901]{L04429 Arthur und Louise Schnitzler an Gustav Schwarzkopf, 17. 4. 1901}
\nopagebreak\mylabel{L04429v}
\rehead{ }\normalsize\beginnumbering\briefempfaengerindex{Schwarzkopf, Gustav@\textsc{Schwarzkopf, Gustav}!zzzSchnitzler, Louise@\emph{von Louise Schnitzler}!1901-04-171@{17. 4. 1901}|(be}\briefempfaengerindex{Schwarzkopf, Gustav@\textsc{Schwarzkopf, Gustav}!zzzSchnitzler, Arthur@\emph{von Arthur Schnitzler}!1901-04-171@{17. 4. 1901}|(be}
\toendnotes[C]{\smallbreak\pagebreak[2]}
\correspDesc{Versand  durch Arthur Schnitzler, Louise Schnitzler am 17. 4. 1901 in Bologna
\newline{}Übermittlung  am 18. 4. 1901 in Bologna
\newline{}Zustellung  am 19. 4. 1901 in Wien
\newline{}Erhalt  durch Gustav Schwarzkopf im Zeitraum [18. 4. 1901 – 22. 4. 1901?] in Wien}\toendnotes[C]{\smallbreak}
\Standort{DLA, A:Schnitzler, HS.1985.1.1897.}
\physDesc{Bildpostkarte, 147 Zeichen
\newline{}Handschrift Arthur Schnitzler: Bleistift, deutsche Kurrent
\newline{}Handschrift Louise Schnitzler: Bleistift, deutsche Kurrent
\newline{}Versand: 1) Stempel: »\nobreak{}\oindex{Bologna@\textbf{Bologna}|pwk}Bologna Ferrovia, 18 4–01, \textcolor{gray}{1M}\nobreak{}«.   2) Stempel: »\nobreak{}\oindex{I., Innere Stadt@\textbf{I., Innere Stadt}|pwk}Wien 1/1 1, 19/4. 1901, 8–9½V, Bestellt\nobreak{}«. }\pstart{}{\pb}Austria\oindex{Österreich@\textbf{Österreich}|pw}\pend{}\pstart{}Herrn Guſtav Schwarzkopf\pend{}\pstart{}\textsc{Wien}\oindex{Wien@\textbf{Wien}|pw}\pend{}\pstart{}\textsc{Tiefer Graben 23}\oindex{Wien@\textbf{Wien}!I., Innere Stadt@\textbf{I., Innere Stadt}!Tiefer Graben 23@\textbf{Tiefer Graben 23}|pw}.\pend{}{\bigskip}
\pstart
           \centering{}{\pb}\textcolor{gray}{\textbf{BOLOGNA. PIAZZA VITT. EMANUELE\oindex{Piazza Maggiore [Bologna]@\textbf{Piazza Maggiore [Bologna]}|pw}.}} \pend
           \vspace{1em}
\pstart
           \raggedleft{}{\pb}17. 4. 901.\pend
           \vspace{0.5em}
\pstart
           Herzliche{\\[\baselineskip]}Grüße. \spacefill\mbox{Arthur}\pend
           \leftskip=0em{}\selectlanguage{ngerman}\vspace{1em}
\pstart
           \noindent{}{[}hs. L. Schnitzler:{]} Viele herzliche Grüße\pend
           
\pstart
           Auf baldiges Wiederſehn!{\\[\baselineskip]}\spacefill\mbox{Louise Schnitzler}\pend
           \leftskip=0em{}\selectlanguage{ngerman}\endnumbering\briefempfaengerindex{Schwarzkopf, Gustav@\textsc{Schwarzkopf, Gustav}!zzzSchnitzler, Louise@\emph{von Louise Schnitzler}!1901-04-171@{17. 4. 1901}|)be}\briefempfaengerindex{Schwarzkopf, Gustav@\textsc{Schwarzkopf, Gustav}!zzzSchnitzler, Arthur@\emph{von Arthur Schnitzler}!1901-04-171@{17. 4. 1901}|)be}\mylabel{L04429h}
\begin{anhang}
\end{anhang}\newcommand{\dateiname}{L04429}\newcommand{\titel}{Arthur und Louise Schnitzler an Gustav Schwarzkopf, 17. 4. 1901}\newcommand{\editorInnen}{Herausgegeben von Jahnke, SelmaMüller, Martin Anton}\input{../tex-inputs/latex-leseansicht-abspann}
